\subsection{CONJUGACY OF COMPACT FORMS}

THEOREM 1. Let \(\mathfrak{a}\) be a complex semi-simple Lie algebra.

\begin{enumerate}
\def\labelenumi{\alph{enumi})}
\item
  \(\mathfrak{a}\) has compact (resp. splittable) real forms.
\item
  The group \(\operatorname{Int}(\mathfrak{a})\) operates transitively on the set of compact (resp. splittable) real forms of a.
\end{enumerate}

Let \(\mathfrak{h}\) be a Cartan subalgebra of \(\mathfrak{a}\) . Then \((\mathfrak{a},\mathfrak{h})\) is split (Chap. VIII, §2, no. 1, Remark 2), and has a Chevalley system \((X_{\alpha})\) (Chap. VIII, §4, no. 4, Cor. of Prop. 5). Part \(a)\) now follows from Prop. 2. Let \(\mathfrak{g}\) be a compact real form of \(\mathfrak{a}\) ; we show that there exists \(v\in\mathrm{Int}(\mathfrak{a})\) such that \(v(\mathfrak{a}_{u})=\mathfrak{g}\) . Let \(\mathfrak{t}\) be a Cartan subalgebra of \(\mathfrak{g}\) ; then \(\mathfrak{t}_{(\mathbf{C})}\) is a Cartan subalgebra of \(\mathfrak{a}\) ; since \(\mathrm{Int}(\mathfrak{a})\) operates transitively on the set of Cartan subalgebras of \(\mathfrak{a}\) (Chap. VII, §3, no. 2, Th. 1), we are reduced to the case in which \(\mathfrak{t}_{(\mathbf{C})}=\mathfrak{h}\) . Since \(\mathfrak{g}\) is a compact form, the eigenvalues of the endomorphisms ad \(h\) , for \(h\in\mathfrak{t}\) , are purely imaginary ( \(\S1\) , no. 3, Prop. 1), so the roots \(\alpha\in\mathbb{R}\) map \(\mathfrak{t}\) to \(i\mathbf{R}\) ; this implies that \(\mathfrak{t}=i\mathfrak{h}_{0}\) . Then, by Prop. 3 (no. 2), there exists \(v\in\mathrm{Int}(\mathfrak{a})\) such that \(v(\mathfrak{a}_{u})=\mathfrak{g}\) , hence \(b)\) in the case of compact forms. Finally, let \(\mathfrak{m}_{1}\) and \(\mathfrak{m}_{2}\) be two splittable real forms of \(\mathfrak{a}\) . There exist framings \((\mathfrak{m}_{1},\mathfrak{h}_{1},\mathrm{B}_{1},(X_{\alpha}^{1}))\) and \((\mathfrak{m}_{2},\mathfrak{h}_{2},\mathrm{B}_{2},(X_{\alpha}^{2}))\) (Chap. VIII, §4, no. 1). These extend in an obvious way to bases \(e_{1}\) and \(e_{2}\) of \(\mathfrak{a}\) . An automorphism of \(\mathfrak{a}\) that maps \(e_{1}\) to \(e_{2}\) maps \(\mathfrak{m}_{1}\) to \(\mathfrak{m}_{2}\) ; thus, it suffices to apply Prop. 5 of Chap. VIII, §5, no. 3, to obtain the existence of an element \(u\) of \(\mathrm{Aut}_{0}(\mathfrak{a})=\mathrm{Int}(\mathfrak{a})\) such that \(u(\mathfrak{m}_{1})=\mathfrak{m}_{2}\) .

Remark. We shall see much later a general classification of real forms of a complex semi-simple Lie algebra.

COROLLARY 1. Let \(\mathfrak{g}\) and \(\mathfrak{g}'\) be two compact real Lie algebras. Then \(\mathfrak{g}\) and \(\mathfrak{g}'\) are isomorphic if and only if the complex Lie algebras \(\mathfrak{g}_{(\mathbf{C})}\) and \(\mathfrak{g}'_{(\mathbf{C})}\) are isomorphic.

The condition is clearly necessary. Conversely, assume that \(\mathfrak{g}_{(\mathbf{C})}\) and \(\mathfrak{g}'_{(\mathbf{C})}\) are isomorphic. Let c (resp. \(c'\) ) be the centre of g (resp. \(g'\) ) and s (resp. \(s'\) ) the derived algebra of g (resp. \(g'\) ). Then \(\mathfrak{c}_{(\mathbf{C})}\) and \(\mathfrak{c}'_{(\mathbf{C})}\) are the centres of \(\mathfrak{g}_{(\mathbf{C})}\) and \(\mathfrak{g}'_{(\mathbf{C})}\) , respectively, and hence are isomorphic; it follows that the commutative algebras c and \(c'\) are isomorphic. Similarly, \(\mathfrak{s}_{(\mathbf{C})}\) and \(\mathfrak{s}'_{(\mathbf{C})}\) are isomorphic, hence s and \(s'\) , which are compact real forms of two isomorphic complex semi-simple Lie algebras, are isomorphic by Th. 1 b).

COROLLARY 2. Let \(\mathfrak{a}\) be a complex Lie algebra. The following conditions are equivalent:

\begin{enumerate}
\def\labelenumi{(\roman{enumi})}
\item
  \(\mathfrak{a}\) is reductive.
\item
  There exists a compact real Lie algebra \(\mathfrak{g}\) such that \(\mathfrak{a}\) is isomorphic to \(\mathfrak{g}(\mathbf{C})\) .
\item
  There exists a compact Lie group \(\mathbf{G}\) such that \(\mathfrak{a}\) is isomorphic to \(\mathrm{L}(\mathrm{G})_{(\mathbf{C})}\) .
\end{enumerate}

By Def. 1 of §1, no. 3, conditions (ii) and (iii) are equivalent and imply (i). If \(\mathfrak{a}\) is reductive, it is the direct product of a commutative algebra, which clearly has a compact real form, and a semi-simple algebra which has one by Th. 1 \(a\) ), hence (i) implies (ii).

COROLLARY 3. Let \(\mathfrak{a}_1\) and \(\mathfrak{a}_2\) be two complex semi-simple Lie algebras. The compact real forms of \(\mathfrak{a}_1 \times \mathfrak{a}_2\) are the products \(\mathfrak{g}_1 \times \mathfrak{g}_2\) , where, for \(i = 1, 2\) , \(\mathfrak{g}_i\) is a compact real form of \(\mathfrak{a}_i\) .

Indeed, there exists a compact real form \(g_{1}\) (resp. \(g_{2}\) ) of \(a_{1}\) (resp. \(a_{2}\) ); then \(g_{1} \times g_{2}\) is a compact real form of \(a_{1} \times a_{2}\) . The corollary now follows from Th. 1 b), applied to \(a_{1}, a_{2}\) and \(a_{1} \times a_{2}\) .

Note that it follows from Cor. 3 above that a compact real Lie algebra g is simple if and only if the complex Lie algebra \(\mathfrak{g}_{(\mathbf{C})}\) is simple. We say that g is of type \(A_{n}\) , or \(B_{n}, \ldots\) , if \(\mathfrak{g}_{(\mathbf{C})}\) is of type \(A_{n}\) , or \(B_{n}, \ldots\) (Chap. VIII, §2, no. 2). By Cor. 1 above, two compact simple real Lie algebras are isomorphic if and only if they are of the same type.

Let G be an almost simple connected compact Lie group (Chap. III, §9, no. 8, Def. 3). We say that G is of type \(A_{n}\) , or \(B_{n}\) , \ldots, if its Lie algebra is of type \(A_{n}\) , or \(B_{n}\) , \ldots. Two simply-connected almost simple compact Lie groups are isomorphic if and only if they are of the same type.

